\section{Introduction}
\label{sec:introduction}

The Einstein vacuum equations determine spacetime from initial geometry
only as long as the evolution remains a Cauchy development. Four\`es-Bruhat
proved local existence and geometric uniqueness \cite{ChoquetBruhat1952};
Choquet-Bruhat and Geroch established a maximal globally hyperbolic
development, unique up to an isometry preserving the initial data
\cite{ChoquetBruhatGeroch1969}. An extension beyond that development need
not be determined by the same data. Strong cosmic censorship asks whether
such extensions are excluded for generic initial data in an appropriate
regularity class. Penrose's discussion of the strong formulation relates
this question to global hyperbolicity and to the infinite blueshift at
an inner horizon \cite[Section~12.3.2]{Penrose1979}.

Kerr's rotating vacuum solutions \cite{Kerr1963} make this question
particularly sharp. Their subextremal interiors contain Cauchy horizons,
so the exact solutions admit continuations beyond their maximal Cauchy
developments \cite{Carter1968}. The conjectural mechanism is instability under perturbation.
Here we prove a local, Baire-generic $C^1$ form of future strong cosmic
censorship near each fixed rotating subextremal Kerr bridge. The data are
smooth, complete, two-ended vacuum data with no symmetry condition.
Genericity is measured in a weighted smooth topology; the neighborhood
requires smallness of only its tenth seminorm. The precise theorem and
the definition of a future timelike exit appear in
Section~\ref{sec:phase}. The neighborhood depends on the Kerr center,
and the conclusion concerns every such exit of the full maximal
globally hyperbolic development, allowing nonvacuum ambient metrics.

\subsection{Why extension regularity matters}

The physical analyses of Poisson and Israel \cite{PoissonIsrael1990}
and Ori \cite{Ori1991,Ori1992} describe how blueshift can produce
mass inflation and a weak tidal singularity at an inner horizon.
Rigorous nonlinear work in the spherically symmetric
Einstein--Maxwell--scalar-field model established that continuous
extendibility can coexist with instability of curvature and mass
\cite{Dafermos2003,Dafermos2005}. In that model, Luk and Oh proved a
generic $C^2$ inextendibility theorem for admissible complete two-ended
asymptotically flat data, combining exterior lower bounds with an
interior instability argument \cite{LukOh2019I,LukOh2019II}.
Building on their work, Sbierski \cite[Theorem~4.31]{Sbierski2022}
proved future locally Lipschitz inextendibility for generic sufficiently
small spherically symmetric perturbations of two-ended subextremal
Reissner--Nordstr\"om data in the same matter model, allowing extensions
that do not preserve spherical symmetry.
The matter model, spherical symmetry and topology in those results
are part of their statements.

For the vacuum equations without symmetry, Luk \cite{Luk2018}
constructed stable local weak null singularities from singular
characteristic data: in the constructed coordinates the metric extends
continuously while its Christoffel symbols fail to be locally square
integrable. Dafermos and Luk
\cite{DafermosLuk2025} prove persistence of a continuously extendible
Cauchy-horizon portion from suitable Cauchy data already inside a
black hole. Hintz \cite[Theorem~13.1]{Hintz2026} proves nonlinear
exterior stability throughout the subextremal range for small data
with specified finite asymptotic power--log expansions and a faster
decaying remainder. Combined with the interior theorem of Dafermos
and Luk, this gives $C^0$ Cauchy-horizon stability for that data class.
Its asymptotic assumptions differ from the symbol-bounded tails
allowed in our weighted phase space.

Gurriaran \cite[Theorem~4.5.1]{Gurriaran2026} obtains locally Lipschitz
inextendibility near timelike infinity from interior spacelike data
with prescribed nonlinear Price asymptotics and a nonvanishing
nonaxisymmetric leading curvature amplitude. Luk and Sbierski
\cite[Theorem~1.2]{LukSbierski2026} obtain continuous extendibility
and locally Lipschitz inextendibility across a specified Cauchy-horizon
portion from characteristic event-horizon data satisfying quantitative
upper and lower radiation bounds. These results explain the interior
instability mechanisms and the importance of the incoming radiation.
The theorem here addresses a different data question: a residual
subset of a neighborhood of complete vacuum Cauchy data, without
assuming a radiation lower bound, and with every future exit covered.

A $C^1$ metric has continuous connection coefficients but need not
have bounded classical curvature. Curvature growth alone therefore does
not obstruct a $C^1$ extension; our argument uses parallel transport.
Sbierski's holonomy method \cite{Sbierski2022} supplies a direct
geometric precedent; his later criterion
\cite[Theorem~3.14]{Sbierski2026} excludes locally Lipschitz extensions
across the specified weak-null boundary under
stable signed integrated-curvature blowup and the specified geometric
hypotheses. Our proof constructs finite loops for which the transport
bound forced by a $C^1$ extension fails. It proves the required
quantitative estimates within the complete Cauchy evolution.

\subsection{The foundation and the two-sign construction}

The companion \cite{Foundation} supplies the complete near-Kerr
background evolution, exterior constraint inverse, and constrained
linear packet construction. We state the particular geometric and
linear conclusions in Sections~\ref{test:section} and
\ref{sec:background}, including their data hypotheses, derivative
orders and parameter dependence. The proofs below use those
intermediate results; the companion's final $C^2$ inextendibility
conclusion is not a prerequisite.

The companion \cite[Theorem~1.2]{KerrWeakConnection} excludes continuous
nondegenerate future extensions with locally square-integrable connection
for a local residual set near each fixed rotating subextremal Kerr bridge.
That stronger inextendibility conclusion builds on the intermediate packet,
constraint, and geometric results developed here, together with its own
nonlinear comparison and averaging over families of loops. The present
paper proves the pointwise transport obstruction at $C^1$ regularity;
the square-integrable-connection theorem is not an input to its proof.

The constraint step belongs to the development of localized vacuum
deformations initiated by Corvino, Corvino--Schoen and
Chru\'sciel--Delay \cite{Corvino2000,CorvinoSchoen2006,ChruscielDelay2003}.
Here an inner ball must remain exactly unchanged while a decaying
exterior correction is retained. This makes the early evolution
identical and leaves room for a late experiment. Complex Gaussian
beams and their superpositions allow an oscillatory packet to cross
real caustics \cite{Ralston1982,Tanushev2008,LiuRunborgTanushev2013}.
The present construction returns that packet to the original initial
surface and controls the Einstein wave-gauge constraints,
matching of normal jets and an interval whose length grows with the
observation depth; those estimates are proved in the companion and
in Section~\ref{sec:background}.

The analytic difficulty is that a packet capable of producing the
needed signal need not dominate its own nonlinear comparison error.
We prepare two exact vacuum perturbations with opposite linear parts,
using the same constraint inverse and the same background gauge.
Their pure-packet quadratic forcing agrees. Subtracting the two
evolution equations cancels that forcing and gives an estimate for
the difference of the two errors. When the curvature observations
are then subtracted, the background term cancels and the remaining
quadratic curvature terms are small enough to control. One of the
two observations must therefore be large. Section~\ref{prep:section}
establishes the signed cancellation at the entry cylinder, and
Section~\ref{energy:section} propagates it through the interior.
The energy estimate keeps one common growth rate at every fixed
derivative order and on every intermediate time interval; this is
what allows the subtraction to preserve its gain even when an
individual error exceeds the linear packet.

\subsection{From one future exit to one finite contradiction}

The logical starting point is a hypothetical future $C^1$ exit.
In a regular extension chart, bounded connection coefficients give
a linear bound for the holonomy of loops of small developed length.
Here developed length means the length of the loop's velocity measured
in a frame transported along that same loop. Section~\ref{test:section}
shows that one bound holds along an open family of finite-duration
timelike geodesics. These requirements form countably many closed
tests on the initial data. It is enough to destroy each test by an
arbitrarily small perturbation in any prescribed finite seminorm.

The comparison development is an auxiliary Cauchy development that
contains the entire initial bridge and an interior region described
by two increasing optical coordinates. Its two branch regimes have
one coordinate tending to infinity and the other to a finite limit;
in the corner regime both coordinates diverge. These are coordinate
limits, with no boundary points adjoined to the perturbed spacetime.
A passing test supplies a geodesic approaching one of the two
noncorner branches of the comparison development. The reason is
geometric: the geodesic flow preserves canonical volume on transverse
sections of the unit mass shell, while corner tracks eventually occupy
strips of arbitrarily small such volume. An open family of initial
seeds cannot consist entirely of corner tracks. This argument does
not impose a bound on curvature along the selected observer.

At a late finite point of that branch, the signed packet construction
produces the curvature separation just described. Section~\ref{attach:section}
attaches both experiments to developments of the entire corrected
initial surface, retaining their asymptotic tails. It also chooses
nearby observing seeds that reach the prescribed observation point.
The comparison solves the perturbed unit mass shell for an absolute
canonical momentum; it avoids dividing an error by the vanishing
optical lapse.

Finally, Section~\ref{holo:section} rescales a null frame adapted to
that observer. Exact null identities control the potentially growing
connection entries. The resulting loop family has exceptionally short
developed length and controlled finite derivatives of its holonomy
with respect to loop size. A finite interpolation formula extracts
the quadratic holonomy coefficient, which is curvature. The assumed
linear transport bound makes this coefficient smaller than the signed
signal. The relevant scales can already be seen here. At optical depth
$T$, the lapse satisfies $a_p\asymp e^{-\kappa T}$ for a fixed
$\kappa>0$. We choose packet amplitude $A_T=e^{-7\kappa T/8}$
and loop parameter $h_T=e^{-\kappa T/16}$. The signed construction
gives a scaled curvature component larger than $A_T$, whereas the
transport test and finite interpolation would give the bound
\[
 e^{o(T)}(a_p/h_T+h_T^{16})
 =e^{o(T)}\bigl(O(e^{-15\kappa T/16})+e^{-\kappa T}\bigr)
 =o(A_T).
\]
The factor $e^{o(T)}$ represents losses smaller than every fixed
positive exponential rate, at the finite derivative orders in use.
This strict separation is the finite contradiction that destroys the test.
Section~\ref{sec:completion} then applies the Baire theorem and
checks the order of all choices. Figure~\ref{fig:proof-map-signed-tests}
records this route through the proof.

\begin{figure}[!htbp]
\centering
\begin{tikzpicture}[
 font=\small,
 box/.style={draw=black!65,rounded corners=2pt,fill=black!3,
  align=center,text width=4.55cm,minimum height=2.5cm,inner sep=6pt},
 arrow/.style={-{Stealth[length=5pt]},thick,draw=black!75}]
\node[box] (exit) at (0,0) {\textbf{Any future $C^1$ exit}\\
 an open family of observers\\has bounded short-loop transport\\
 {\footnotesize Section~\ref{test:section}}};
\node[box] (branch) at (5.4,0) {\textbf{A noncorner observer}\\
 canonical phase volume\\excludes an all-corner family\\
 {\footnotesize Section~\ref{test:section}}};
\node[box] (pair) at (10.8,0) {\textbf{Paired exact perturbations}\\
 opposite linear packets;\\one common propagation rate\\
 {\footnotesize Sections~\ref{prep:section}--\ref{energy:section}}};
\draw[arrow] (exit.east)--(branch.west);
\draw[arrow] (branch.east)--(pair.west);
\node[box] (category) at (0,-3.35) {\textbf{Baire genericity}\\
 every closed test is destroyed\\in every data neighborhood\\
 {\footnotesize Section~\ref{sec:completion}}};
\node[box] (loop) at (5.4,-3.35) {\textbf{A short-loop contradiction}\\
 finite holonomy jets recover\\the scaled curvature coefficient\\
 {\footnotesize Section~\ref{holo:section}}};
\node[box] (signal) at (10.8,-3.35) {\textbf{One large signed signal}\\
 the background cancels;\\the odd error and quadratic\\curvature terms are controlled\\
 {\footnotesize Sections~\ref{attach:section}--\ref{holo:section}}};
\draw[arrow] (pair.south)--(signal.north);
\draw[arrow] (signal.west)--(loop.east);
\draw[arrow] (loop.west)--(category.east);
\end{tikzpicture}
\caption{The route from a possible future exit to a finite contradiction.
Each experiment is attached to the full development of its complete
corrected datum. The noncorner branch locates an observation; the final
contradiction concerns parallel transport in the full development.}
\label{fig:proof-map-signed-tests}
\end{figure}

\FloatBarrier
